% Transcription — fonds Alexandre Grothendieck, shelfmark 79, batch 4 (pages 61-80).
% Established from the facsimile archives/batches/79/batch-04.pdf (a mirror of
% https://grothendieck.umontpellier.fr/79.pdf), per the skill
% transcribe-grothendieck.
%
% Pass: Opus 5.5 (claude-opus-5-5), 2026-10-03 — first pass, unchecked against the pages by a human.
%
% All twenty pages are transcribed; they carry the author's pagination 30-39.
% Page 70 holds only three lines and a struck heading.
%
% Pages 61-70 — §IV: for a free module M of rank 2 over a ring Λ, the set
%   S_M of basis vectors mod ±1, the edges A_M, the faces F_M (triples with
%   x+y+z=0), the "type" maps τ to Π_M = (det M)^*/±1, the group
%   G_M = Gl(M)/±1 and its action; the two cases 2·1_Λ = 0 and ≠ 0; whether
%   A_M determines F_M (no, unless the types are given).
% Pages 71-80 — §V « Récapitulation »: the structures \tilde P_M, \tilde P_{M,i}
%   on M^{**}, oriented faces, antipodism; a lemma and corollary (automorphisms
%   fixing a frame are trivial when Λ = Λ_0); theorem for Λ = Z or Z/nZ:
%   Gl(M) ≅ Aut(\tilde P_M); a « Scholie » with two groupoid equivalences;
%   the action of Λ^* by homotheties. Breaks off mid-sentence.
\documentclass[11pt,a4paper]{article}
\input{../preamble/grothendieck}

\folder{79}
\batch{4}
\pages{61}{80}
\dating{[à partir de 1980]}
\watermark{Édition de démonstration}
\foldertitle{2-polyèdres réguliers triangulés, et modules libres de rangs 2 sur les anneaux $\Lambda=\mathbb{Z}/n\mathbb{Z}$ et $\Lambda=\mathbb{Z}$ : notes manuscrites (s.d.).}

\begin{document}

\note{le lot commence au milieu d'un texte paginé par l'auteur : la numérotation des formules (156 et suivantes) et les renvois (82), (118) p.~23 désignent des pages antérieures au lot.}

\page{61}
\note{p. 30 de l'auteur. La numérotation des formules, à partir de (156), continue celle des pages qui précèdent ce lot ; le paragraphe IV s'ouvre ici.}
\note{notations fixées pour ce lot : $\mathfrak{P}_2(E)$, $\mathfrak{P}_3(E)$ rendent le P barré qu'il emploie pour l'ensemble des parties à deux (resp. trois) éléments de $E$ ; $\mathrm{Rep}(M)$ est l'ensemble des repères (bases ordonnées) $(x,y)$ de $M$ ; $\dot{x}$ est l'image de $x$ dans $S_M$.}

\textbf{IV}\quad Soient \struck{\ill{}} $\Lambda$ un anneau non nul, $M$ un module libre de rang $2$ sur $\Lambda$,
\[
\text{(156)}\quad
\left\{
\begin{array}{l}
\widetilde{S}_M = \{x\in M \mid \exists\, y\in M,\ (x,y)\in \mathrm{Rep}(M)\}\\[4pt]
S_M = \widetilde{S}_M/(\pm\mathrm{id}_M)
\end{array}
\right.
\]
\note{devant la seconde ligne, un numéro (157) est biffé.}
de telle façon que l'on ait canoniquement
\[
\text{(157)}\qquad
\widetilde{S}_M \longrightarrow S_M
\]
revêtement de degré constant, égal à $2$ si $2\cdot 1_\Lambda\neq 0$, égal à $1$ si $2\cdot 1_\Lambda=0$, i.e. $\Lambda$ anneau $\mathbb{F}_2$-alg.

Soient
\[
\text{(158)}\quad
\left\{
\begin{array}{l}
\text{\struck{$\widetilde{A}_M = \mathrm{Rep}(M)/\text{symétrie }(x,y)\mapsto(y,x)$}}\\
\text{\struck{$\phantom{\widetilde{A}_M} = \{\{x,y\}\in\mathfrak{P}_2(M)\mid \{x,y\}\text{ base de }M\}$}}\\[4pt]
A_M = \text{\struck{$\mathrm{Rep}(M)/\pm\mathrm{id}_M$}}\ \text{\struck{$\widetilde{A}_M/\pm\mathrm{id}_M$}}\\
\phantom{A_M ={}} \{\{s,t\}\in\mathfrak{P}_2(S_M)\mid \exists\,(x,y)\in\mathrm{Rep}(M),\\
\phantom{A_M ={}\{\{s,t\}\in\mathfrak{P}_2(S_M)\mid{}} \{s,t\}=\{\dot{x},\dot{y}\}\}
\end{array}
\right.
\]
\note{la définition de $\widetilde{A}_M$ est encadrée et biffée ; seule celle de $A_M$ subsiste.}

\marginal{NB Il y a 4 choix (et non 2 !) de $x,y$ tels que $(s,t)=(\dot{x},\dot{y})$, donc 8 choix tels que $\{s,t\}=\{\dot{x},\dot{y}\}$.}\note{« 4 » cerclé ; le début de la note, en marge gauche, porte des mots biffés.}
\struck{j'ai donc \ill{}}

\struck{(159)\quad $\widetilde{A}_M\longrightarrow A_M$ de degré constant, égal à $2$ si $2\cdot 1_\Lambda\neq 0$, à $1$ si $2\cdot1_\Lambda=0$.}

\note{tout le bas de la page, depuis la ligne précédente, est biffé de longs traits obliques ; il se lit néanmoins.}
\struck{On a une inclusion canonique}
\[
\text{\struck{(160)\quad $\widetilde{A}_M \hookrightarrow \mathfrak{P}_2(S_M)$}}
\]
\struck{du fait que si $\{x,y\}$ est une base de $M$, alors $y\notin\{x,-x\}$ donc $\{\dot{x},\dot{y}\}\in\mathfrak{P}_2(S_M)$,} \struck{et bien sûr $\{\dot{x},\dot{y}\}=\{-x,$} \struck{et bien sûr \ill{} \ill{}. Posant $\{x,y\}^{\cdot}$, on a}
\[
\text{\struck{$\bigl(-\{x,y\}\bigr)^{\cdot} = \{x,y\}^{\cdot}$}}
\qquad\text{\struck{où $-\{x,y\} \overset{\mathrm{déf}}{=} \{-x,-y\}$.}}
\]

\page{62}
On considère $A_M$ comme l'ens. des \emph{arêtes} d'une structure de graphe sur $S_M$. \struck{Le groupe} On a une application naturelle
\[
\text{(161)}\qquad A_M \xrightarrow{\ \tau\ \text{ou}\ \tau_A\ } \Pi_M \overset{\mathrm{déf}}{=} (\det M)^{*}/\pm 1_\Lambda
\]
\marginal{on regardera $\Pi_M$ comme un \uncertain{quotient} de $A_M$, grâce à $\tau$.}
déduite par passage au quotient de l'application
\[
\text{(162)}\quad
\left\{
\begin{array}{rcl}
\widetilde{A}_M & \xrightarrow{\ \widetilde{\tau}\ } & \Pi_M\\
\{x,y\} & \longmapsto & x\wedge y \ \text{mod \emph{signe}}
\end{array}
\right.
\]
\marginal{NB. ne change pas si on échange $x$ et $y$, \ill{} \ill{} $x\wedge y$ \ill{} des signes.}

Posons, pour $i\in\Pi_M$,
\[
\text{(163)}\qquad A_{M,i} = \tau^{-1}(\{i\}),
\]
qu'on considère comme l'ens. des arêtes d'une structure de graphe \struck{\ill{}} $\Sigma(i)$ sur $S_M$.

Le groupe $\mathrm{Gl}_\Lambda(M)$ opère sur $(S_M,A_M)$ par transport de structure, ainsi que sur $\Pi_M$, et $\tau$ (161) commute aux opérations de $\mathrm{Gl}_\Lambda(M)$. Le sous-groupe $\{\pm\mathrm{id}_M\}$ de $\mathrm{Gl}_\Lambda(M)$ opère trivialement sur les ens. $S_M$, $\Pi_M$, on

\page{63}
\note{p. 31 de l'auteur.}
conclut que $\mathrm{Gl}_\Lambda(M)$ opère par l'intermédiaire d'une \struck{opération}
\[
\text{(164)}\qquad \text{opération de } G_M = \mathrm{Gl}_\Lambda(M)' = \mathrm{Gl}_\Lambda(M)/\pm 1 \text{ sur } S_M,
\]
respectant la structure $A_M\subset\mathfrak{P}_2(S_M)$, et l'ens. quotient $\Pi_M$ de $A_M$, \uncertain{et} pour $i\in\Pi$, $g\in G_M$, $g_{S_M}$ est un isom. de la structure $(S_M,A_{M,i})$ avec la structure $(S_M, A_{M,g_{\Pi_M} i})$. En particulier, si on \struck{\ill{}} considère le composé
\[
\text{(165)}\qquad
G_M \xrightarrow{\ \det\ } \Lambda^{*} \xrightarrow{\ \mathrm{can}\ } \Lambda^{*}/\pm 1_\Lambda \overset{\mathrm{déf}}{=} Z'_\Lambda ,
\]
\note{sous la flèche composée $G_M\to Z'_\Lambda$, un arc porte le nom $\delta_M$.}
\marginal{\emph{NB} l'opération de $G_M$ sur $\Pi_M$ respecte aussi la structure de $Z'_\Lambda$-torseur, où $Z'_\Lambda=\Lambda^{*}/\pm1_\Lambda$, et $G_M$ opère sur $\Pi_M$ via $\delta_M$.}
le stabilisateur dans $G_M$ \struck{de tout} d'un point de $\Pi_M$ \struck{est} est le sous-groupe $G^{\circ}_M = \operatorname{Ker}\delta_M$ intervenant dans la suite exacte
\[
\text{(166)}\qquad
1 \longrightarrow G^{\circ}_M \longrightarrow G_M \xrightarrow{\ \delta_M\ } Z'_\Lambda \longrightarrow 1 ,
\]
où $G^{\circ}_M = \mathrm{Sl}_\Lambda(M)/\pm 1$ et $G_M = \mathrm{Gl}'_\Lambda(M)$,\note{ces deux égalités sont écrites sous la suite, rattachées à ses termes par des signes $=$ verticaux ; la première se lit « $\mathrm{Sl}'_\Lambda(M)\pm1$ ».}
et il en résulte que c'est aussi le sous-groupe de $G_M$ qui stabilise la structure de \struck{graphe} $\Sigma(i)$, pour un $i\in\Pi_M$ quelconque donné.

\page{64}
On voit que $G^{\circ}_M$ est transitif sur \struck{toutes} les ens. $A_{M,i}$ ($i\in\Pi_M$), et \struck{simplement} transitif sur l'ens. $\overrightarrow{A}_{M,i}$ des arêtes orientées du graphe \struck{$(S_M,A_{M,i})$} $\Sigma(i)$ ; la transitivité provient du fait que si $(x,y),(x',y')\in\mathrm{Rep}(M)$ sont tels que $x\wedge y=\pm x'\wedge y'$, alors l'unique $u\in\mathrm{Gl}_\Lambda(M)$ tel que $(x',y')=u\cdot(x,y)$ est dans $\mathrm{Gl}_\Lambda(M)^{\pm1}$, donc $u\in G^{\circ}_M$. [La \emph{simple} transitivité est fausse, sauf si $2\cdot1_\Lambda=0$ donc $S_M\simeq M^{**}$\note{lecture douteuse de l'exposant.}, $A_M$ = ens. des bases non ordonnées $\{x,y\}$ de $M$, \struck{\ill{}}, cf.\ plus bas] \struck{provient du fait que \ill{} $u\cdot(x,y)\in\{\pm(x,y)\}$ implique $u\in\{\pm\mathrm{id}_M\}$ i.e. $u=1_{G_M}$. Il en résulte que l'on ait}

Il en résulte que le groupe $G_M$ est transitif sur $A_M$, et \struck{simplement} transitif sur l'ens. $\overrightarrow{A}_M$.

Introduisons aussi \struck{\ill{}}
\[
\text{(167)}\quad
\begin{array}{l}
F_M = \bigl\{ f\in\mathfrak{P}_3(S_M) \bigm| \exists\, x,y,z\in M\ \bigl(x+y+z=0,\ x\wedge y\in(\det M)^{*}\bigr)\\
\phantom{F_M = \bigl\{ f\in\mathfrak{P}_3(S_M) \bigm|{}} \text{tels que } f=\{\dot{x},\dot{y},\dot{z}\}\bigr\}
\end{array}
\]
\note{numéro lu (157) sur la page ; la suite de la numérotation demande (167). Devant $M$ dans « $x,y,z\in M$ », une lettre surchargée, peut-être $\widetilde{S}$, reste incertaine ; avant $x+y+z=0$ un mot est biffé.}
Alors
\[
\text{(168)}\quad
\left\{
\begin{array}{l}
\forall f\in F_M,\ \forall a\in\mathfrak{P}_2(f), \text{ on a } a\in A_M ;\\
\text{de plus, } \tau(a) \text{ ne dépend que de } f, \text{ soit } \tau(f)\in\Pi_M
\end{array}
\right.
\]
\note{à gauche de cette accolade, une première rédaction est biffée et illisible.}

\page{65}
\note{p. 32 de l'auteur.}
donc on a une application
\[
\text{(169)}\qquad F_M \xrightarrow{\ \tau_F\ } \Pi .
\]

Par ailleurs, on a à distinguer deux cas
\[
\text{(170)}\quad \text{A)}\ \ 2\cdot1_\Lambda=0 \quad\text{alors}\quad \forall a\in A_M\ \exists !\, f\in F_M,\ f\supset a
\]
\note{il écrit « $f\ni F_M$ » pour $f\in F_M$, ici et dans (171).}

\begin{tikzcd}[column sep=large]
A_M \arrow[r, "\text{d'où}"] \arrow[rr, bend right=30, "\tau_A"'] & F_M \arrow[r, "\tau_F"] & \Pi
\end{tikzcd}

\note{numéro (170 bis). Au-dessus de la première flèche, « déf. » biffé puis « d'où » ; dessous : « \uncertain{hom.} surjectif de degré 3 » ; sous la seconde, un mot, \uncertain{homom.}, que le dessin ne place pas sûrement.}

et le cas
\[
\text{(171)}\quad \text{B)}\ \ 2\cdot1_\Lambda\neq0 \quad\text{alors } \forall a\in A_M, \text{ il existe exactement \emph{deux} } f\in F_M \text{ tels que } f\supset a .
\]

Dans le cas A), on a
\[
\text{(172)}\qquad \overrightarrow{A}_M \xleftarrow{\ \sim\ } \underbrace{\mathrm{Rep}(M)/\pm\mathrm{id}_M}_{\mathrm{Rep}'(M)} \qquad (\text{car } 2\cdot1_\Lambda=0)
\]
\marginal{donc $G_M$ est \add{\uncertain{tel}} simplement transitif sur $\overrightarrow{A}_M$, i.e. $\overrightarrow{A}_M$ est un $G_M$-torseur.}

Dans le cas B), l'application canonique
\[
\text{(173)}\quad
\left\{
\begin{array}{rcl}
\mathrm{Rep}'(M) & \longrightarrow & \overrightarrow{A}_M\\
(x,y)^{\cdot} & \longmapsto & (\dot{x},\dot{y})
\end{array}
\right.
\qquad \text{est de degré deux, } (2\cdot1_\Lambda\neq0)
\]
elle se factorise en
\[
\text{(174)}\qquad
\mathrm{Rep}'(M) \xrightarrow{\ \sim\ } \mathrm{Rep}(P_M) \overset{\mathrm{déf}}{=} \bigl\{(s,t,u)\in S_M \bigm| \{s,t,u\}\in F_M\bigr\}
\]
\[
\text{\struck{$\mathrm{Rep}(M)$}}\quad (x,y)^{\cdot} \longmapsto \bigl(\dot{x},\dot{y},(-x-y)^{\cdot}=(x+y)^{\cdot}\bigr)
\]
suivie de
\[
\text{(175)}\quad
\left\{
\begin{array}{rcl}
\mathrm{Rep}(P_M) & \longrightarrow & \overrightarrow{A}_M\\
(s,t,u) & \longmapsto & (s,t)
\end{array}
\right.
\quad\Big|\ \text{qui est \add{bien} de degré 2, par (171).}
\]

\page{66}
Dans le cas A, on appellera \emph{repère du} graphe $P_M = (S_M, A_M)$ un élément de $\overrightarrow{A}_M$, i.e.
\[
\text{(176)}\qquad \mathrm{Rep}(P_M) \overset{\mathrm{déf}}{=} \overrightarrow{A}_M \qquad (\text{cas } 2\cdot1_\Lambda=0),
\]
dans le cas B), on définit $\mathrm{Rep}(P_M)$ \struck{\ill{}} par (174). Dans tous les cas, l'ens. des repères \uncertain{abs.} $P_M$ est un torseur sous $G_M$.

\emph{Relations entre les données} $A_M$, $F_M$. Je voudrais m'assurer que ces données sur $S_M$
\[
\text{(177)}\qquad
\left[\begin{array}{l} A_M\subset\mathfrak{P}_2(S_M)\\ A_M\to\Pi\end{array}\right],
\quad
\left[\begin{array}{l} F_M\subset\mathfrak{P}_3(S_M)\\ F_M\to\Pi\end{array}\right]
\]
se déterminent mutuellement, comme en B),\note{lecture incertaine de « comme en B) », écrit sous la ligne.} En tous cas, dans le cas A) \struck{$(2\cdot1_\Lambda\neq0)$}, on a
\marginal{$A_M\to\Pi$, $F_M\to\Pi$ sont \uncertain{interprétées} ici comme des partitions de $A_M$, $F_M$ \ill{} \ill{} $\Pi$ \ill{} \ill{}.}
\[
\text{(177)}\qquad A_M = \{a\in\mathfrak{P}_2(S_M)\mid \exists\, f\in F_M,\ f\supset a\}
\]
\note{le numéro (177) est écrit deux fois sur la page.}
(168, 171)\add{(170)}, et d'autre part je voudrais, \add{dans le cas B $(2\cdot1_\Lambda\neq0)$},
\[
\qquad F_M = \{f\in\mathfrak{P}_3(S_M)\mid \forall a\in\mathfrak{P}_2(f), \text{ on a } a\in A_M\}
\]
\marginal{relation fausse !!}
On a bien $F_M\subset\{\ \}$, par (168), inversement \struck{a-t-on} a-t-on $F_M\supset\{\ \}$ ?

C'est \uncertain{ultra} \emph{faux}, même sur $\Lambda=\mathbb{Z}/n\mathbb{Z}$ ou $\Lambda=\mathbb{Z}$, cela signifiant que si $x,y,z\in M$ sont tels que $x\wedge y$, $y\wedge z$, $z\wedge x$ sont dans $\det M^{*}$, alors on peut, \ill{} \uncertain{trouver}

\page{67}
\note{p. 33 de l'auteur.}
$\varepsilon,\varepsilon',\varepsilon''\in\{\pm1\}$ tels que l'on ait $\varepsilon x+\varepsilon' y+\varepsilon'' z=0$, ou encore qu'on a une égalité de la forme
\[
(*)\qquad z = \varepsilon x+\varepsilon' y \qquad (\varepsilon,\varepsilon'\in\{\pm1_\Lambda\})
\]
ce qui est évidemment faux. Par contre, \struck{pour \ill{} si $i\in$} si on se donne, en plus de $A_M$, l'application $\tau_A : A_M\to\Pi_M$, i.e. la décomposition de $A_M$ suivant les types $i\in\Pi$, on reconstitue $F_M$ comme
\[
\text{(178)}\quad
\left\{
\begin{array}{l}
F_M = \coprod_{i\in\Pi} F_{M,i}\\[4pt]
F_{M,i} = \{f\in\mathfrak{P}_3(S_M)\mid \forall a\in\mathfrak{P}_2(f), \text{ on a } a\in A_{M,i}\}
\end{array}
\right.
\]
cette dernière relation en effet signifie que si $x,y,z\in M$ sont tels que
\[
x\wedge y,\ y\wedge z,\ z\wedge x \in i=\{e,-e\}\subset(\det M)^{*}
\]
alors on doit avoir une relation $(*)$ ci-dessus, en effet. \struck{\ill{}} Si $(x,y)\in\mathrm{Rep}(M)$, on aura $z=\varepsilon x+\varepsilon' y$, $\varepsilon,\varepsilon'\in\Lambda$, et on écrira, en multipliant \add{$x$} cette relation par
\[
\varepsilon' x\wedge y = x\wedge z \qquad \text{d'où } \varepsilon'\in\{\pm1_\Lambda\}
\]
et de même $\varepsilon\in\{\pm1_\Lambda\}$, cqfd.

Dans le cas A) par contre $(2\cdot1_\Lambda=0)$,

\page{68}
il semble\add{rait} que la structure
\[
\left\{
\begin{array}{l}
F_M\subset\mathfrak{P}_3(S_M)\\
F_M\xrightarrow{\ \tau_F\ }\Pi_M
\end{array}
\right.
\]
soit nettement plus fine que la structure $A_M\subset\mathfrak{P}_2(S_M)$, $A_M\xrightarrow{\tau_A}\Pi_M$ qui s'en déduit, via (177) (168) : \struck{Une application de $M$ \ill{}}

Si $\Lambda$ \uncertain{est} un corps $k$ de car.\ $2$, \struck{\ill{}} corps premier $\mathbb{F}_2$, alors une \struck{bijection} \add{permutation} de $M^{*}=M\smallsetminus\{0\}$ (où $M\simeq k^2$) peut fort bien transformer couples de pts \struck{qui sont une base} \add{colinéaires} \struck{alignés} en couples \add{colinéaires} \struck{alignés} --- donc non-\uncertain{colinéaires} \uncertain{en non-colinéaires} --- \uncertain{bases en bases} --- et, je suppose, de façon compatible avec les \uncertain{jauges}, i.e.\ en transformant \struck{\ill{}} deux bases de \uncertain{même jauge} en \struck{\ill{}} bases de \uncertain{même jauge}, sans pour autant respecter la relation $x+y+z=0$, pour $\bigl(x\wedge y=y\wedge z=z\wedge x\in(\det M)^{*}\bigr)$ …
\note{la lecture « jauge », pour la classe de $x\wedge y$, est conjecturale ; le mot revient trois fois sous la même forme.}
\marginal{\emph{à élucider} \ill{} --- \uncertain{éventuellement}, en utilisant \uncertain{la} structure d'addition \ill{} \uncertain{et} \ill{} \ill{} de $\Pi_M=(\det M)^{*}$, et exigeant qu'elle soit respectée par les \uncertain{permutations}, \ill{} on arrive à déduire \ill{} $F_M$ de $A_M$, $\tau_A : A_M\to\Pi_M$ …}

À vrai dire, la structure que nous étudions \add{alors}, en car.\ $2$, se réduisant à l'ens.\ de $\widetilde{S}_M = M^{**}$ (ens.\ des éléments de $M$ faisant partie d'une base), l'ens.\ $\widetilde{F}_M \subset \mathfrak{P}_3(\widetilde{S})$\note{le signe d'inclusion paraît barré.} formé de parties $\{x,y,z\}$ de $\widetilde{S}_M$ satisfaisant $x+y+z=0$, $x\wedge y\ (=y\wedge z=z\wedge x)\in(\det M)^{*}$, et de l'application
\[
\widetilde{\tau} : \widetilde{F}_M \longrightarrow (\det M)^{*} = \widetilde{\Pi},
\qquad \{x,y,z\}\longmapsto x\wedge y = y\wedge z = z\wedge x
\]
(178),\note{renvoi à (178) écrit au pied de la page, à gauche.}

\page{69}
\note{p. 34 de l'auteur.}
peut avoir un intérêt sur un anneau quelconque $\Lambda$, pas néc.\ de car.\ $2$ i.e.\ tel que $2\cdot1_\Lambda=0$, à cela près que l'application $\widetilde{\tau}$ est alors à valeurs non dans $\widetilde{\Pi}$, mais dans $\Pi=\widetilde{\Pi}/\{\pm1_\Lambda\}$
\[
\text{(178 bis)}\quad
\left\{
\begin{array}{rcl}
\widetilde{\tau} : \widetilde{F}_M & \longrightarrow & \Pi = (\det M^{*})/\{\pm1\}\\
\{x,y,z\} & \longmapsto & x\wedge y \ \mathrm{mod}\ \pm1
\end{array}
\right.
\]
--- si on veut une application à valeurs dans $\det M^{*}$, il faudra prendre
\[
\text{(179)}\quad
\left\{
\begin{array}{rcl}
\widetilde{\tau}_{\mathrm{or}} : \underset{\text{« faces orientées »}}{\widetilde{F}.\mathrm{or}._M} & \longrightarrow & \widetilde{\Pi} = \det M^{*}\\[10pt]
(x,y,z),\ (y,z,x),\ (z,x,y) & \longmapsto & x\wedge y = y\wedge z = z\wedge x
\end{array}
\right.
\]
Notons qu'on a \add{même} une application analogue sur $F.\mathrm{or}._M \simeq \widetilde{F}.\mathrm{or}._M/\pm\mathrm{id}_M$, puisque l'on a
\[
\widetilde{\tau}_{\mathrm{or}}\bigl(\vec{f}\,\bigr)\ \text{\struck{$(x,y,z)$}} = \widetilde{\tau}_{\mathrm{or}}\bigl(-\mathrm{id}_M(\vec{f}\,)\bigr)
\]
pour une face orientée de \uncertain{$\widetilde{P}_M$} --- d'où
\[
\text{(180)}\qquad F.\mathrm{or}._M \xrightarrow{\ \tau_{\mathrm{or}}\ } \text{\struck{\ill{}}}\ \widetilde{\Pi}
\]
qui \struck{\ill{}} s'insère dans le diag.\ commutatif

\begin{tikzcd}
F.\mathrm{or}._M \arrow[r] \arrow[d, "\mathrm{can}"'] & \widetilde{\Pi} \arrow[d]\\
F_M \arrow[r] & \Pi
\end{tikzcd}

\note{numéro (181).}

\page{70}
et la question se pose dans quelle mesure cette structure supplémentaire se déduit des autres structures….

\struck{\uncertain{Récapitulation} \ill{}}
\note{titre souligné puis biffé ; le reste de la page est blanc.}

\page{71}
\note{p. 35 de l'auteur.}

\textbf{V}\quad \emph{Récapitulation}

Soit $\Lambda$ un anneau commut.\ non nul, $M$ un $\Lambda$-module libre de rang $2$. On lui associe deux « structures triangulées » (ens.\ de sommets $S$, plus un \add{ens.\ de} parties \add{$F$} \struck{\ill{}} $\subset\mathfrak{P}_3(S)$ \struck{\ill{}} ternaires dans $S$), dont la deuxième n'est intéressante que si $2\cdot1_\Lambda\neq0$ (autrement, elles coïncident avec la première). Les structures triangulées \struck{sont} en question sont « partitionnées » (i.e.\ $F\subset\mathfrak{P}_3(S)$ est muni d'une partition, via une application dans un ens.\ d'indices convenable $I$) --- donc la structure envisagée peut être considérée comme un ens.\ de structures triangulées sur le même ens.\ $S$, indexé par $I$. Les deux structures envisagées dépendent fonctoriellement de $M$, et la deuxième est le quotient de la première par $\{\pm\mathrm{id}_M\}\subset\mathrm{Gl}_\Lambda(M)$.

A)
\[
\text{(1)}\quad
\left\{
\begin{array}{l}
\widetilde{S}_M = M^{**} = \{x\in M\mid \exists\, y\in M,\ (x,y)\in\mathrm{Rep}(M)\}\\[4pt]
\widetilde{F}_M = \Bigl\{\{x,y,z\}\in\mathfrak{P}_3(\widetilde{S}_M) \Bigm| \begin{array}{l} x+y+z=0\\ x\wedge y\ (=y\wedge z=z\wedge x)\in(\det M)^{*}\end{array}\Bigr\}
\end{array}
\right.
\]
\[
\text{(2)}\qquad \overrightarrow{\Pi}_M = \det M^{*},\quad \Pi_M = (\det M^{*})/\{\pm1_\Lambda\}
\]

\page{72}
On pose
\[
\text{(3)}\qquad
\underset{\substack{\text{ens. des \emph{faces}}\\ \text{\emph{orientées} de } \widetilde{P}_M=(\widetilde{S}_M,\widetilde{F}_M)}}{\overrightarrow{\widetilde{F}}_M}
= \Bigl\{(f,\omega)\Bigm| f\in\widetilde{F}_M,\ \omega\in\underset{\substack{\text{ens. des ordres}\\ \text{circ. sur } f}}{\underline{\omega}_f}\Bigr\}
\]
\note{après le $F$ de $\overrightarrow{\widetilde{F}}_M$, deux lettres sont noircies, ici et en (5).}
\[
\text{(4)}\qquad \vec{f}\longmapsto\check{\vec{f}}
\qquad\text{donné par}\qquad (f,\omega)\longmapsto(f,\omega')
\]
(NB.\ on \struck{\ill{}} risque une notation $-\vec{f}$ à la face orientée déduite de $\vec{f}$ par transport de structure par $-\mathrm{id}_M$)
\[
\text{(5)}\quad
\left\{
\begin{array}{l}
\overrightarrow{\widetilde{F}}_M \xrightarrow{\ \vec{\widetilde{\tau}}_{\mathrm{or}}\ } \overrightarrow{\Pi}_M\\[4pt]
\widetilde{\tau}\bigl(\vec{f}_{x,y,z}\bigr) = x\wedge y\ \ (=y\wedge z=z\wedge x)
\end{array}
\right.
\]
\marginal{on note $\vec{f}_{x,y,z}\in\widetilde{F}.\mathrm{or}._M$ la face orientée ($\{x,y,z\}$, $x\to y\to z\to x$)}
On a
\[
\vec{\widetilde{\tau}}_{\mathrm{or}}\bigl(\check{\vec{f}}\,\bigr) = -\,\vec{\widetilde{\tau}}_{\mathrm{or}}\bigl(\vec{f}\,\bigr)
\]
d'où par passage au quotient une appl.
\[
\text{(6)}\qquad \widetilde{F}_M \xrightarrow{\ \widetilde{\tau}\ } \Pi_M .
\]
Pour $i\in\Pi$, on pose
\[
\text{(7)}\quad
\left\{
\begin{array}{ll}
\widetilde{F}_{M,i} = \widetilde{\tau}^{-1}(\{i\}) & \text{d'où } \widetilde{F}_M = \coprod_{i\in\Pi}\widetilde{F}_{M,i}\\[4pt]
\widetilde{P}_{M,i} & \text{est le « polyèdre triangulé » } (\widetilde{S}_M,\widetilde{F}_{M,i})
\end{array}
\right.
\]
\struck{\uncertain{Proposition}} On pose
\[
\text{(8)}\qquad
\underset{\text{\emph{arêtes} de } \widetilde{P}_M}{\widetilde{A}_M} = \{a\in\mathfrak{P}_2(\widetilde{S}_M)\mid \exists\, f\in\widetilde{F}_M,\ f\supset a\}
\]
\[
\Bigl[=\{\{x,y\}\in\mathfrak{P}_2(\widetilde{S}_M)\mid (x,y)\in\mathrm{Rep}(M)\}\Bigr]
\]
et on a

\page{73}
\note{p. 36 de l'auteur.}
(9)\ \struck{\emph{Proposition}} Toute arête de $\widetilde{P}_M$ est contenue dans une \struck{et une seule face} une et une seule face, d'où
\[
\widetilde{A}_M \xrightarrow[\text{degré 3}]{} \widetilde{F}_M
\]
et une application composée

\begin{tikzcd}
\widetilde{A}_M \arrow[r] \arrow[rr, bend right=30, "\widetilde{\tau}_A"'] & \widetilde{F}_M \arrow[r, "\widetilde{\tau}_F"] & \Pi_M
\end{tikzcd}

\note{numéro (10).}
De même, avec la définition évidente des arêtes orientées, on a

\begin{tikzcd}[column sep=large]
\underset{\text{arêtes orientées de } \widetilde{A}_M}{\overrightarrow{\widetilde{A}}_M} \arrow[r, "\text{degré 3}"] \arrow[rr, bend right=25, "\vec{\widetilde{\tau}}_A"'] & \overrightarrow{\widetilde{F}}_M \arrow[r, "\vec{\widetilde{\tau}}_F"] & \overrightarrow{\Pi}_M
\end{tikzcd}

\note{numéro (11).}
NB\quad $\vec{\widetilde{\tau}}_A(\check{\vec{a}}) = -\,\vec{\widetilde{\tau}}_A(\vec{a})$
\marginal{(10), (11) \struck{\ill{}} \ill{} se déduisent \uncertain{par passage aux quotients} \ill{} de \ill{}}

On définit $\widetilde{A}_{M,i}$ \struck{\ill{}} (pour $i\in\Pi_M$) comme l'image inverse de $\{i\}$ par $\widetilde{\tau}_A$, d'où
\[
\text{(12)}\qquad \widetilde{A}_M = \coprod_{i\in\Pi}\widetilde{A}_{M,i} .
\]
\note{la ligne (13) qui suit est encadrée et biffée de hachures.}
\struck{(13)\quad $\widetilde{A}_{M,i}$ et $\widetilde{F}_{M,i}$ se déterminent mutuellement :}
\[
\text{\struck{$\widetilde{A}_{M,i} = \{a\in\mathfrak{P}_2(\widetilde{S}_M)\mid \exists\, f\in\widetilde{F}_{M,i},\ f\supset a\}$}}
\]
\[
\text{\struck{$\widetilde{F}_{M,i} = \{f\in\mathfrak{P}_3(\widetilde{S}_M)\mid \forall a\in\mathfrak{P}_2(f),\ a\in\widetilde{A}_{M,i}\}$ !}}
\]
\struck{\ill{}} On peut considérer que $(\widetilde{F}_{M,i})_{i\in\Pi_M}$ définit une famille de structures triangulées sur le même ens.\ de sommets $\widetilde{S}_M$, paramétrée par $\Pi_M$. Il faudra à l'occasion interpréter la structure
\marginal{\struck{\ill{} une structure}}
\[
\text{(13)}\qquad \Pi_M \text{ torseur sous } Z'_\Lambda = \Lambda^{*}/\{\pm1_\Lambda\},
\]
et aussi
\[
\text{(14)}\qquad \overrightarrow{\Pi}_M \text{ torseur sous } \Lambda^{*} \text{ qui « relève » } \Pi_M .
\]
\marginal{\uncertain{tirer au clair}}\note{les numéros (13) et (14) sont écrits par-dessus d'autres ; un trait vertical relie ces deux lignes à la note marginale.}

\page{74}
\emph{Remarques}\quad Il \emph{n'est pas vrai} que la seule donnée des $\widetilde{A}_{M,i}$ permette de retrouver les $\widetilde{F}_{M,i}$, par la recette naïve --- car \add{pour} \ill{} $f\in\mathfrak{P}_3(\widetilde{S}_M)$, $f=\{x,y,z\}$, il ne suffit pas qu'on ait $a\in\widetilde{A}_{M,i}$ $\forall a\in\mathfrak{P}_2(f)$, i.e.\ $x\wedge y$, $y\wedge z$, $z\wedge x\in i=\{\pm e\}\subset\det M^{*}$, pour avoir $f\in\widetilde{F}_{M,i}$ i.e.\ $x+y+z=0$. En effet, on constate que la condition envisagée signifie simplement que $z=\pm x\pm y$ (et non $z=-(x+y)$). Mais dans le cas $2\cdot1_\Lambda=0$ --- qui est celui où \struck{\ill{}} la structure $\widetilde{P}_M$ \add{\uncertain{elle-même}} \struck{\ill{}} est la plus intéressante --- cela implique bien $x+y+z=0$, i.e.
\[
\text{(15)}\quad \text{si } 2\cdot1_\Lambda=0, \text{ alors }
\widetilde{F}_{M,i} = \{f\in\mathfrak{P}_3(\widetilde{S}_M)\mid \forall a\in\mathfrak{P}_2(f), \text{ on a } a\in\widetilde{A}_{M,i}\}
\]

\emph{Question de connexité des} $\widetilde{P}_{M,i}$. On s'y intéresse, car on obtiendrait qu'un automorphisme de $\widetilde{P}_{M,i}$ (i.e.\ permutation de $\widetilde{S}_M$ qui \struck{\ill{}} \add{transforme} $\widetilde{F}_{M,i}$ en lui-même) qui fixe un \emph{repère} $(x,y,z)$ (i.e.\ un $(x,y,z)\in\widetilde{S}_M^{\,3}$ tel que $\{x,y,z\}\in\widetilde{F}_{M,i}$) doit être l'identité. [Pour pouvoir travailler, il \ill{} \ill{} \ill{} manière que l'automorphisme \ill{} « l'antipodisme » $\in\widetilde{a}_M$ \uncertain{automorphisme} de $\widetilde{P}_{M,i}$ \ill{} induit par $-\mathrm{id}_M$). Alors, si $u$ fixe un $x\in\widetilde{S}_M$,
\marginal{donc il faut considérer \ill{} comme \ill{} \ill{} la structure de $\widetilde{P}_M$ \ill{}}
\note{l'argument continue au-delà de la page 74 ; la ligne marginale, très serrée, n'est lisible que par fragments.}

\page{75}
\note{p. 37 de l'auteur.}
\marginal{\struck{\uncertain{inutile} ?}}
il fixe aussi $-\xi$\note{lettre surchargée : $x$ corrigé en $\xi$, ou l'inverse.}.] S'il fixe $\xi$\,\add{$x$} et $y$ telles que $x\wedge y\in i=\{e,-e\}\subset\overrightarrow{\Pi}=(\det M)^{*}$, alors il fixe aussi \add{$z=$} $-(x+y)$, (\add{donc aussi $x+y$.}) telle nous permet de nous propager : partant d'un \add{$i$-}repère $(x,y)$, \struck{fixé \ill{} les \uncertain{remplaçant} par $(x,y+x)$} i.e.\ à partir d'une $i$-face $(x,y,z)$, en faisant soit \struck{des transpositions} \add{une permutation} sur $(x,y,z)$ --- p.\ ex.\ les deux transpositions entre $x$ et $y$, et entre $y$ et $z$ --- et la transformation $(x,y,z)\longmapsto(x,-y,y-x)$ --- en termes matriciels, un $i$-repère de $M$ $r=(x,y)$, interprété comme iso
\[
r : \Lambda^2 \xrightarrow{\ \sim\ } M
\]
qui soit formé de « composantes » $x,y$ fixes par $u$ (on dit que $r$ est « $u$-fixe »), il en est de même de $r\circ\tau_0^{r}$, $r\circ\tau_1^{r}$, $r\circ\tau_\infty^{r}$ (cf.\ (82)\note{renvoi « p. 15 » écrit au-dessus de (82).}). Cela signifie que si $r$ est $u$-fixe, il en est de même de $r\circ g$, où
\[
g\in\underset{\mathrm{Aut}_\Lambda(\Lambda^2)}{\underline{\mathrm{Gl}(2,\Lambda)}} \text{ est dans le sous-groupe engendré par } \tau_0^{r},\tau_1^{r},\tau_\infty^{r} .
\]
Or ce sous-groupe est l'image, par
\[
\mathrm{Sl}(2,\mathbb{Z}) \longrightarrow \mathrm{Sl}(2,\Lambda)\subset\mathrm{Gl}(2,\Lambda)
\]
du sous-groupe de $\mathrm{Sl}(2,\mathbb{Z})$ engendré par les matrices $T_i^{r}$ (transpositions comme ci-dessus à

\page{76}
coeff.\ entiers). Comme on sait que ce sous-groupe est $\mathrm{Sl}(2,\mathbb{Z})$ tout entier, et que si $\Lambda_0=\mathrm{Im}(\mathbb{Z}\to\Lambda)$, \struck{\ill{}} l'homom.
\[
\mathrm{Sl}(2,\mathbb{Z}) \longrightarrow \underset{\substack{\cap\\ \mathrm{Sl}(2,\Lambda)\subset\mathrm{Gl}(2,\Lambda)}}{\mathrm{Sl}(2,\Lambda_0)} \quad\text{est surjectif}
\]
on voit que le sous-groupe de $\mathrm{Sl}(2,\Lambda)$ engendré par les $T_i^{r}$ n'est autre que \struck{$\mathrm{Sl}(2,\Lambda_0)$.}
\[
\underset{\substack{\cap\\ \mathrm{Sl}(2,\Lambda)}}{\mathrm{Sl}(2,\Lambda_0)} = \Bigl\{\begin{pmatrix} a & b\\ c & d\end{pmatrix} \Bigm| \begin{array}{l} a,b,c,d\in\Lambda_0=\mathrm{Im}(\mathbb{Z}\to\Lambda)\\ ad-bc=1\end{array}\Bigr\},
\]
en fonction opérant \uncertain{ci-dessus} \struck{\uncertain{à droite}} sur l'ens.\ des repères, via
\[
(x,y)\longmapsto(ax+by,\ cx+dy)
\]
Ainsi, on a prouvé

\emph{Lemme}\quad Soit $u$ un automorphisme de $P_{M,i}$, qui fixe \add{deux éléments $x,y\in\widetilde{S}_M$ formant} une $i$-base $(x,y)$ de $M$ (i.e.\ $x\wedge y\in i\subset\overrightarrow{\Pi}_M=\det M^{*}$ ---), \struck{Alors} supposons de plus que $u$ commute à l'antipodisme $-\mathrm{id}_M$. Alors $u$ fixe \struck{les} tous les éléments de $\widetilde{S}_M$ qui sont de la forme $ax+by$, avec $a,b\in\Lambda_0=\mathrm{Im}(\mathbb{Z}\to\Lambda)$, et premiers entre eux (i.e.\ $a\Lambda+b\Lambda=\Lambda$ i.e.\ $\exists\, c,d\in\Lambda$ avec $ad-bc=1$).\note{il écrit « $ad-cc=1$ ».}

\page{77}
\note{p. 38 de l'auteur.}
\emph{Corollaire}\quad Supposons $\Lambda=\Lambda_0$ i.e.\ $\Lambda$ est un anneau « \uncertain{simple} ». Alors tout automorphisme de $P_{M,i}$ qui fixe un repère (i.e.\ les sommets d'une face) et commute à l'antipodisme est l'identité.\note{le mot entre guillemets, peut-être « absolu », reste douteux.}

On trouve donc

\emph{Théorème}\quad Soit $\Lambda\simeq\mathbb{Z}$ ou $\mathbb{Z}/n\mathbb{Z}$, \add{$M$ un $\Lambda$-module libre de rang $2$.} Alors
\begin{itemize}
\item[a)] Pour une structure $\tfrac12$-symplectique $i\in\underset{\Pi_M}{\underbrace{(\det M^{*})/\pm1}}$, l'\uncertain{homom.} \struck{du groupe des automorphismes}
\[
\text{(16)}\qquad \mathrm{Gl}_\Lambda(M)^{\pm} \longrightarrow \text{\struck{\ill{}}}\ \mathrm{Aut}(\widetilde{P}_{M,i})^{\text{antipod}}
\]
\struck{est injectif (\uncertain{c'est évident}) et} est un isomorphisme [car $\mathrm{Gl}_\Lambda(M)^{\pm}$ est simplement transitif sur l'ens.\ des bases $\tfrac12$-symplectiques relatives à $i$].
\note{sous $\mathrm{Aut}(\widetilde{P}_{M,i})^{\text{antipod}}$, une accolade porte : « automorphismes qui commutent à l'antipodisme ».}
\item[b)] \struck{\ill{}} \add{Considérons sur} $\widetilde{S}_M=M^{**}$ la structure $\widetilde{P}_M$ définie par la famille $(\widetilde{F}_{M,i})_{i\in\Pi_M}$ de structures triangulées, et l'antipodisme, qui est un automorphisme commun à toutes ces structures triangulées. Alors
\[
\text{(17)}\qquad \mathrm{Gl}_\Lambda(M) \xrightarrow{\ \sim\ } \mathrm{Aut}(\widetilde{P}_M) .
\]
\end{itemize}
\marginal{il suffit \uncertain{de regarder} la structure $\widetilde{P}_M$ définie par $\widetilde{F}_M$ tout seul, \uncertain{sans} \uncertain{partition}}

\page{78}
\emph{Scholie}\quad \add{Soit $n\in\mathbb{N}\smallsetminus\{1\}$, d'où $\Lambda_n\simeq\mathbb{Z}$ si $n=0$, $\mathbb{Z}/n\mathbb{Z}$ sinon.} Par les principes généraux, il s'ensuit qu'on \struck{a\uncertain{} réciproquement} \add{\uncertain{deux}} équivalences de groupoïdes
\[
(M,i)\longmapsto\widetilde{P}_{M,i}
\]
\[
\text{(18)}\quad
\underset{\substack{\text{\struck{Groupoïde} des modules libres}\\ \text{$\tfrac12$-symplectiques de rang 2 sur } \Lambda_n}}{\phantom{x}}
\xrightarrow{\ \approx\ }
\underset{\substack{\text{groupoïde des structures triangulées}\\ \text{avec antipodisme (automorphisme involutif)}\\ \text{isomorphes à une structure type}\\ \text{convenable (dépendant \struck{de \ill{}} de $n$)}}}{\phantom{x}}
\]
\[
M\longmapsto\widetilde{P}_M
\]
\[
\text{(19)}\quad
\underset{\substack{\text{groupoïde des modules libres de}\\ \text{rang 2 sur } \Lambda_n\ \text{\struck{($n\in\mathbb{N}$, …)}}}}{\phantom{x}}
\xrightarrow{\ \approx\ }
\underset{\substack{\text{groupoïde des structures } (S,F,\Pi,\tau,a)\\ \text{où } F\subset\mathfrak{P}_3(S) \text{ une str. triangulée sur } S,\\ \text{décomposée en famille \add{disjointe} indexée par } \Pi\\ \text{(i.e. on a } \tau:F\to\Pi \text{ surj.),}\\ a \text{ est un automorphisme involutif}\\ \text{(antipodisme) de la structure précédente}}}{\phantom{x}}
\]
\note{dans (18) et (19), les deux termes de chaque équivalence sont écrits en prose sur la page ; on les a placés sous un support vide.}
\marginal{NB Si $2\cdot1_\Lambda=0$ (i.e., ici, $\Lambda=\mathbb{F}_2$) il est inutile d'introduire l'antipodisme, qui est l'identité. Si $2\cdot1_\Lambda\neq0$ \struck{\ill{}} l'antipodisme est \uncertain{déjà} \ill{} dans \ill{} \uncertain{ses points fixes}…}

Il conviendrait de caractériser axiomatiquement, par des propriétés intrinsèques, les structures qui figurent au deuxième membre. De plus, il faudrait comprendre, en termes de ces dernières, comment reconstruire $M$, et par là aussi,
\marginal{cf.\ suggestion (23) ci-dessous}
\note{un trait vertical double borde, dans la marge gauche, ce dernier paragraphe.}

\page{79}
\note{p. 39 de l'auteur.}
les structures supplémentaires reliant les données $S$, $F$ etc.\ à $\Lambda$ (cf.\ (13) (14) (5) et aussi (118) p.~23).

Enfin, dans le cas où on ne suppose plus $\Lambda=\Lambda_0$, on voudrait introduire sur $\widetilde{P}_{M,i}$ \add{\uncertain{resp.} sur $\widetilde{P}_M$} un élément de structure supplémentaire (peut-être l'opération de $\Lambda^{*}$ sur $\widetilde{S}_M$) qui permette néanmoins d'interpréter $\mathrm{Gl}_\Lambda(M)^{\pm}$ resp.\ $\mathrm{Gl}_\Lambda(M)$ comme le groupe d'automorphismes de cette structure.
\marginal{il y a aussi l'opération de $\Lambda^{*}$ \struck{\ill{}} sur $\widetilde{S}_M$}
\note{ces deux paragraphes sont bordés dans la marge gauche d'un double trait vertical.}

\note{le paragraphe B) qui suit est barré de traits obliques.}
\struck{\uncertain{Si} $2\cdot1_\Lambda\neq0$} \add{On suppose dorénavant $2\cdot1_\Lambda\neq0$.}
\struck{B) Le groupe $\{\pm\mathrm{id}_M\}$ opère librement sur $\widetilde{S}_M$, $\widetilde{A}_M$, $\widetilde{F}_M$ [NB si $\{x,y,z\}\in\widetilde{F}_M$, alors les éléments $x,y,z,-x,-y,-z$ sont tous distincts…]. On pose}

\struck{Revenons à} \add{\uncertain{donc} \uncertain{pour} $\Lambda$ quelconque} $\widetilde{P}_M$, et considérons l'opération de $\Lambda^{*}$ sur $\widetilde{S}_M$, par les $\widetilde{T}_\lambda$
\[
\text{(20)}\qquad \widetilde{T}_\lambda x = \lambda\cdot x
\]
alors $\Lambda^{*}$ respecte toutes ces structures ($\widetilde{F}_M\subset\mathfrak{P}_3(\widetilde{S}_M)$ et la partition de $\widetilde{F}_M$ en $\widetilde{F}_{M,i}$) --- l'opération de $\mathrm{Gl}_\Lambda(M)$ sur $\widetilde{S}_M$ et commute à $\Lambda^{*}\simeq\mathrm{Centre}(\mathrm{Gl}_\Lambda(M))$, en respectant aussi
\marginal{NB La donnée de l'antipodisme \uncertain{entre} \uncertain{dans} cette opération de $\Lambda^{*}$}
\note{la phrase se poursuit à la page suivante.}

\page{80}
les structures $(\widetilde{F}_{M,i})_{i\in\Pi}$.

Soit $u$ un automorphisme de cette structure, commutant aux \struck{\ill{}} $\widetilde{T}_\lambda$, et fixant un repère $(x,y,z)$, i.e.\ fixant les éléments d'une base $(x,y)$. \struck{Alors} Notons que pour \uncertain{tout} élément $\xi\in\widetilde{S}_M=M^{**}$, s'il fixe $\xi$, il fixe les $T_\lambda\xi=\lambda\xi$, $\lambda\in\Lambda^{*}$, en particulier aussi $-\xi$ --- \struck{\ill{}} \add{Comme il} \struck{par le Lemme p.~37', qui pour $a,b\in$} \add{\uncertain{fixe}} \add{si} $x$ fixe $x$ et les $\lambda y$ ($\lambda\in\Lambda^{*}$), on voit qu'il \struck{induit l'identité} fixe les $\tau_A(\{x,\lambda y\})=\lambda i$ (où $i=\tau_A(\{x,y\})$), donc il induit l'identité sur $\Pi_M$, i.e.\ \struck{\ill{}} \add{c'est un} automorphisme pour toutes les structures $\widetilde{F}_{M,i}$.
\marginal{il suffit qu'il fixe $x$, $y$}
\note{la rédaction de cette phrase est très surchargée ; « p. 37' » renvoie au Lemme de la page 76 du fonds.}

\struck{En vertu du Lemme p.~37' avec $x,y$, il fixe aussi les $ax+by$ avec …} Considérons les matrices, $\lambda\in\Lambda^{*}$
\[
\left\{
\begin{array}{l}
\underline{\lambda}_0(\lambda) = \begin{pmatrix} 1 & -\lambda\\ 0 & 1\end{pmatrix}\\[10pt]
\lambda_1 = \begin{pmatrix} 0 & 1\\ -1 & 0\end{pmatrix}\\[10pt]
\lambda_\infty = \begin{pmatrix} 1 & -1\\ 1 & 0\end{pmatrix}
\end{array}
\right.
\]
\note{les noms des trois matrices se lisent $\lambda_0(\lambda)$, $\lambda_1$, $\lambda_\infty$ ; on attendrait plutôt des $\tau$, comme à la page 75.}
Alors l'ens.\ des bases $(x,y)$ « fixes sous $u$ » est stable par composition avec les matrices précédentes et leurs inverses, donc aussi les composés de telles. Si $\Lambda^{*}$ engendre $\Lambda$ \add{additivement} (p.ex.\ $\Lambda$ $\tfrac12$local) alors \struck{\ill{}} l'on
\note{la page s'arrête au milieu de la phrase ; l'argument continue au-delà du lot.}

\end{document}
